% !TeX root = ../main.tex
% 第 5 章 可比较的网测度 (对应原书 pp. 64--74)

\chapter{可比较的网测度}\label{chap:5}

\section{网测度的构造}

\newterm{comparable-net-measure}{可比较网测度}{comparable net measure}是研究 Hausdorff 测度的极为有用的工具.\newterm{net-measure}{网测度}{net measure}的性质比 Hausdorff 测度方便得多，但仍可构造得对许多用途而言与之等价.网测度最先由 \textcite{besicovitch1952} 使用，用以证明 $\mathcal{H}^s$-测度无穷的闭集包含测度为正但有限的子集；后来 \textcite{marstrand1954b} 在研究集合的 Descartes 积的 Hausdorff 测度时使用了它.本章我们特别关注这两个应用.

这里我们把注意力限于 Euclid 空间上的一种基本形式的网测度.更广泛的讨论见 \textcite{rogers1970} 的 2.7~节，那里描述了 Davies 与 Rogers 的工作.

网测度的构造方式与 Hausdorff 测度类似，只是在定义中使用一个受限的覆盖集合类 $\mathcal{N}$ 而非全体集合.类 $\mathcal{N}$ 被选为一个集合的``网''，其本质性质是：若 $U_1,U_2\in\mathcal{N}$，则 $U_1\cap U_2=\emptyset$、$U_1\subset U_2$、$U_2\subset U_1$ 三者必居其一.此外，我们假定 $\mathcal{N}$ 的每个集合只含于有限多个其他集合之中.特别地，给定 $\mathcal{N}$ 中集合的任意族，可以通过去掉那些含于其他集合之中的集合，找到具有同样并集的不交子族.正如我们将看到的，可以构造与 Hausdorff 测度``可比较''的网测度，即两个测度之比上下都有界的网测度.

为构造 $\mathbb{R}^n$ 上的 $s$ 维网测度 $\mathcal{M}^s$，令 $\mathcal{N}$ 为所有 $n$ 维\newterm{binary-cubes}{半开二进方体}{binary cubes, half open}之族，即形如
\[
  [2^{-k}m_1,2^{-k}(m_1+1))\times[2^{-k}m_2,2^{-k}(m_2+1))\times\cdots\times[2^{-k}m_n,2^{-k}(m_n+1))
\]
的集合，其中 $k$ 为非负整数，$m_1,\ldots,m_n$ 为整数.（若 $n=1$ 或 $2$，网 $\mathcal{N}$ 由半开二进区间或正方形组成.）设 $E\subset\mathbb{R}^n$，$\delta>0$，定义
\begin{equation}\label{eq:5.1}
  \mathcal{M}^{s}_{\delta}(E)=\inf\sum_{i=1}^{\infty}|S_i|^{s},
\end{equation}
其中下确界取遍 $E$ 的所有由 $\mathcal{N}$ 中集合 $\{S_i\}$ 构成的可数 $\delta$-覆盖.由网的性质，只需考虑 $E$ 的由 $\mathcal{N}$ 中两两不交的集合构成的覆盖.于是 $\mathcal{M}^s_{\delta}(E)$ 是 $\mathbb{R}^n$ 上在 $\mathbb{R}^n$ 的有界子集上取有限值的外测度.令
\[
  \mathcal{M}^{s}(E)=\lim_{\delta\to 0}\mathcal{M}^{s}_{\delta}(E)=\sup_{\delta>0}\mathcal{M}^{s}_{\delta}(E),
\]
便得到 $\mathbb{R}^n$ 上的度量外测度 $\mathcal{M}^s$.由\cref{thm:1.5}，Borel 集是 $\mathcal{M}^s$-可测的（事实上 Souslin 集亦然）.与\cref{thm:1.6} 的证明相仿可证 $\mathcal{M}^s$ 是正则外测度.

我们来证明 $\mathcal{M}^s$ 与 $\mathcal{H}^s$ 确实是可比较的测度；这本质上是 \textcite{besicovitch1952} 的工作.

\begin{theorem}\label{thm:5.1}
存在只依赖于维数 $n$ 的常数 $b_n$，使得对每个 $E\subset\mathbb{R}^n$，当 $0<\delta<1$ 时
\begin{equation}\label{eq:5.2}
  \mathcal{H}^{s}_{\delta}(E)\leq\mathcal{M}^{s}_{\delta}(E)\leq b_n\mathcal{H}^{s}_{\delta}(E),
\end{equation}
且
\begin{equation}\label{eq:5.3}
  \mathcal{H}^{s}(E)\leq\mathcal{M}^{s}(E)\leq b_n\mathcal{H}^{s}(E).
\end{equation}
\end{theorem}

\begin{proof}
由外测度的定义立得 $\mathcal{H}^s_{\delta}(E)\leq\mathcal{M}^s_{\delta}(E)$.（Hausdorff 测度的下确界取遍更大的覆盖集合类.）

设 $U$ 为满足 $0<|U|\leq\delta$ 的任意集合，令 $k$ 为满足
\begin{equation}\label{eq:5.4}
  2^{-k-1}\leq|U|<2^{-k}
\end{equation}
的整数，并设 $S$ 为与 $U$ 相交的边长 $2^{-k}$ 的二进方体.则 $U$ 含于由 $S$ 及其直接邻域组成的 $3^n$ 个边长 $2^{-k}$、直径 $2^{-k}n^{1/2}$ 的二进方体之族.把每个这样的方体再剖分为 $2^{n^2}$ 个更小的方体，则 $U$ 含于 $b_n=3^n2^{n^2}$ 个直径为
\begin{equation}\label{eq:5.5}
  2^{-k}n^{1/2}2^{-n}\leq 2^{1-n}n^{1/2}|U|\leq|U|\leq\delta
\end{equation}
的二进方体之中，这里用了 \eqref{eq:5.4}.现在设 $\{U_i\}$ 是 $E$ 的由任意集合构成的 $\delta$-覆盖.对每个 $i$ 有 $U_i\subset\bigcup_{j=1}^{b_n}S_{ij}$，其中 $\{S_{ij}\}_{j=1}^{b_n}$ 是由 $b_n$ 个直径至多为 $|U_i|\leq\delta$ 的方体构成的族.

于是 $E\subset\bigcup_{i}\bigcup_{j}S_{ij}$，且
\[
  \sum_{i=1}^{\infty}\sum_{j=1}^{b_n}|S_{ij}|^{s}\leq b_n\sum_{i=1}^{\infty}|U_i|^{s},
\]
故 \eqref{eq:5.2} 右端的不等式由 $\mathcal{M}^s_{\delta}$ 与 $\mathcal{H}^s_{\delta}$ 的定义得出.

令 $\delta\to 0$ 即得不等式 \eqref{eq:5.3}.
\end{proof}

本节最后证明两条关于集合的单调序列的技术性引理，本章后面将会用到.

\begin{lemma}\label{lem:5.2}
设 $\{E_j\}$ 是 $\mathbb{R}^n$ 的紧子集的递减序列.则对任意 $\delta>0$，
\begin{equation}\label{eq:5.6}
  2^{s}\mathcal{H}^{s}_{\delta}\Bigl(\lim_{j\to\infty}E_j\Bigr)\geq\lim_{j\to\infty}\mathcal{H}^{s}_{2\delta}(E_j).
\end{equation}
\end{lemma}

\begin{proof}
设 $\{U_i\}$ 是 $\lim E_j$ 的任意 $\delta$-覆盖.对每个 $i$，令 $V_i$ 为包含 $U_i$ 的开集，$|V_i|\leq 2|U_i|$，并令 $V$ 表示开集 $\bigcup_{i}V_i$.我们断言对某个整数 $j$ 有 $E_j\subset V$.否则 $\{E_j\setminus V\}$ 是非空紧集的递减序列，由紧性的一个初等推论，它有非空的极限集（$\lim E_j\setminus V$）.于是
\[
  2^{s}\sum|U_i|^{s}\geq\sum|V_i|^{s}\geq\mathcal{H}^{s}_{2\delta}(E_j)\geq\lim_{j\to\infty}\mathcal{H}^{s}_{2\delta}(E_j),
\]
考虑 $\lim E_j$ 的所有此类 $\delta$-覆盖即得 \eqref{eq:5.6}.
\end{proof}

若 $\{E_j\}$ 是 $\mathbb{R}^n$ 的子集的任意递增序列，则由\cref{lem:1.3} 有 $\lim_{j\to\infty}\mathcal{M}^s(E_j)=\mathcal{M}^s\bigl(\lim_{j\to\infty}E_j\bigr)$.然而这一结论对外测度 $\mathcal{M}^s_{\delta}$ 也成立.它是\newterm{increasing-sets-lemma}{``递增集引理''}{increasing sets lemma}的一个特殊情形；该引理在 Hausdorff 测度理论中相当重要，\textcite{rogers1970}（2.6~节）与 \textcite{davies1970} 及其所引文献中有详细讨论.我们需要如下特殊情形.

\begin{lemma}\label{lem:5.3}
设 $\{E_j\}$ 是 $\mathbb{R}^n$ 的子集的递增序列，且每个 $E_j$ 都是二进方体的有限并.则
\[
  \mathcal{M}^{s}_{\delta}\Bigl(\lim_{j\to\infty}E_j\Bigr)=\lim_{j\to\infty}\mathcal{M}^{s}_{\delta}(E_j).
\]
\end{lemma}

\begin{proof}
记 $E=\lim_{j\to\infty}E_j=\bigcup_jE_j$.由于对一切 $j$ 有 $\mathcal{M}^s_{\delta}(E_j)\leq\mathcal{M}^s_{\delta}(E)$，只需在右端有限的假设下证明 $\mathcal{M}^s_{\delta}(E)\leq\lim_{j\to\infty}\mathcal{M}^s_{\delta}(E_j)$.由于 $E_j$ 是二进方体的有限并，$\mathcal{M}^s_{\delta}(E_j)$ 定义中的下确界可由某个由二进方体构成的有限不交族取得.设对每个 $j$，$\mathcal{S}_j$ 是这样一个 $E_j$ 的有限 $\delta$-覆盖，使
\begin{equation}\label{eq:5.7}
  \sum_{C\in\mathcal{S}_j}|C|^{s}=\mathcal{M}^{s}_{\delta}(E_j),
\end{equation}
并设 $\mathcal{S}_j$ 是具有该性质的数值上最小的方体族之一.若 $S\in\mathcal{S}_j$，则 $S$ 必含 $E_j$ 的一点.该点也属于 $E_{j+1}$，从而属于某个 $T\in\mathcal{S}_{j+1}$.由网的性质，要么 $S\subset T$，要么 $T\subset S$.若 $T$ 是 $S$ 的真子集，则要么把 $S$ 换成 $\mathcal{S}_{j+1}$ 中含于 $S$ 的那些方体以缩小 $\sum_{C\in\mathcal{S}_j}|C|^s$，要么把 $\mathcal{S}_{j+1}$ 中含于 $S$ 的方体换成单个方体 $S$，以缩小 $\sum_{C\in\mathcal{S}_{j+1}}|C|^s$ 或该和式的项数.这两种情形都与 $\mathcal{S}_j$ 或 $\mathcal{S}_{j+1}$ 的极小性矛盾.于是我们得出：若 $S\in\mathcal{S}_j$，则存在 $T\in\mathcal{S}_{j+1}$ 使 $S\subset T$.

设 $\{C_1,C_2,\ldots\}$ 是从 $\bigcup_{j=1}^{\infty}\mathcal{S}_j$ 中去掉所有含于该族其他方体之中的方体后所得的方体集.则对每个 $j$ 有 $E_j\subset\bigcup_{i=1}^{\infty}C_i$，从而 $E\subset\bigcup_{i=1}^{\infty}C_i$，且
\begin{equation}\label{eq:5.8}
  \mathcal{M}^{s}_{\delta}(E)\leq\sum_{i=1}^{\infty}|C_i|^{s}.
\end{equation}
但由前一段的结论，每个方体 $C_i$ 属于一切充分大的 $j$ 的 $\mathcal{S}_j$.于是给定 $k$，可找到 $j(k)$ 使方体 $C_1,\ldots,C_k$ 都在 $\mathcal{S}_{j(k)}$ 中.利用 \eqref{eq:5.8} 与 \eqref{eq:5.7},
\[
  \mathcal{M}^{s}_{\delta}(E)
  \leq\lim_{k\to\infty}\sum_{i=1}^{k}|C_i|^{s}
  \leq\lim_{k\to\infty}\sum_{C\in\mathcal{S}_{j(k)}}|C|^{s}
  =\lim_{k\to\infty}\mathcal{M}^{s}_{\delta}\bigl(E_{j(k)}\bigr)
  \leq\lim_{j\to\infty}\mathcal{M}^{s}_{\delta}(E_j),
\]
如所需.
\end{proof}

\section{有限测度的子集}

\label{idx:finite-measure-subsets}
\addterm{有限测度}{finite measure, subsets of}{\pageref{idx:finite-measure-subsets}}
本节结果的一个推论是：$s$-集（即具有正的有限 $\mathcal{H}^s$-测度的集合）极为丰富.事实上，给定任何拓扑上像样的、满足 $\mathcal{H}^s(E)>0$ 的集合 $E$，可以找到 $E$ 的紧子集 $F$ 使 $0<\mathcal{H}^s(F)<\infty$.若 $E$ 不是 $\sigma$-有限的 $\mathcal{H}^s$-测度集，这绝非平凡.（$\sigma$-有限集是指可以表示为有限测度集的可数并的集合.）这类结果最先由 \textcite{besicovitch1952} 对闭集得到，并立即被 \textcite{davies1952b} 推广到 Souslin 集.

\begin{theorem}\label{thm:5.4}
设 $E$ 是 $\mathbb{R}^n$ 中满足 $\mathcal{H}^s(E)=\infty$ 的闭子集.

(a) 设 $c$ 为正数.则存在 $E$ 的紧子集 $F$ 使 $\mathcal{H}^s(F)=c$.

(b) 存在 $E$ 的紧子集 $F$ 使 $\mathcal{H}^s(F)>0$，且对某个常数 $b$ 有
\[
  \mathcal{H}^{s}\bigl(B_r(x)\cap F\bigr)\leq br^{s}\qquad(x\in\mathbb{R}^n,\ r\leq 1).
\]
\end{theorem}

\begin{proof}
为使记号相对简单，我们对 $n=1$ 给出证明；高维情形步骤完全相同，只是把二进区间换成二进方体.不妨设 $E$ 有界.由于 $\mathcal{M}^s(E)=\infty$，可找到整数 $m$ 使
\[
  \mathcal{M}^{s}_{2^{-m}}(E)\geq 2^{s}cb_1,
\]
其中 $b_1$ 是\cref{thm:5.1} 中的常数.

归纳地定义 $E$ 的闭子集的递减序列 $\{E_k\}_{m}^{\infty}$.令 $E_m=E$.对 $k\geq m$，通过对每个长度为 $2^{-k}$ 的二进区间指定交来定义 $E_{k+1}$.设 $I$ 是这样的区间.若 $\mathcal{M}^{s}_{2^{-(k+1)}}(E_k\cap I)\leq 2^{-sk}$，就简单地取 $E_{k+1}\cap I=E_k\cap I$.于是
\begin{equation}\label{eq:5.9}
  \mathcal{M}^{s}_{2^{-(k+1)}}\bigl(E_{k+1}\cap I\bigr)=\mathcal{M}^{s}_{2^{-k}}\bigl(E_k\cap I\bigr)\leq 2^{-sk},
\end{equation}
因为在计算 $\mathcal{M}^s_{2^{-k}}$ 时用 $I$ 作为覆盖区间所给出的估计至少不小于用更短区间所得者.另一方面，若 $\mathcal{M}^{s}_{2^{-(k+1)}}(E_k\cap I)>2^{-sk}$，则选取 $E_{k+1}\cap I$ 为 $E_k\cap I$ 的满足 $\mathcal{M}^{s}_{2^{-(k+1)}}(E_{k+1}\cap I)=2^{-sk}$ 的紧子集.这样的子集存在，因为 $\mathcal{M}^{s}_{2^{-(k+1)}}\bigl(E_{k+1}\cap I\cap(-\infty,u]\bigr)$ 关于 $u$ 连续.由于 $\mathcal{M}^{s}_{2^{-k}}(E_k\cap I)=2^{-sk}$，\eqref{eq:5.9} 仍然成立.现在把 \eqref{eq:5.9} 对所有长度为 $2^{-k}$ 的二进区间 $I$ 求和.用于计算 $\mathcal{M}^s_{2^{-k}}$ 与 $\mathcal{M}^s_{2^{-(k+1)}}$ 的覆盖区间必含于某个这样的 $I$ 中，故得
\[
  \mathcal{M}^{s}_{2^{-(k+1)}}(E_{k+1})=\mathcal{M}^{s}_{2^{-k}}(E_k)\qquad(k\geq m).
\]
迭代得
\begin{equation}\label{eq:5.10}
  \mathcal{M}^{s}_{2^{-k}}(E_k)=\mathcal{M}^{s}_{2^{-m}}(E_m)\qquad(k\geq m).
\end{equation}
设 $I$ 是长度为 $2^{-k}$ 的二进区间.若 $m\leq k<r$，则 $E_r\subset E_{k+1}$，故由 \eqref{eq:5.9},
\[
  \mathcal{M}^{s}_{2^{-(k+1)}}(E_r\cap I)\leq\mathcal{M}^{s}_{2^{-(k+1)}}(E_{k+1}\cap I)\leq 2^{-sk}.
\]
于是，在计算 $\mathcal{M}^s_{2^{-k}}(E_r)$ 时，任何长度为 $2^{-k}$ 的区间 $I$ 都可以换成至多 $2^{-(k+1)}$ 长度的区间而不增大下确界的值.故
\[
  \mathcal{M}^{s}_{2^{-(k+1)}}(E_r)=\mathcal{M}^{s}_{2^{-k}}(E_r)\qquad(m\leq k<r),
\]
从而迭代并结合 \eqref{eq:5.10} 得
\begin{equation}\label{eq:5.11}
  \mathcal{M}^{s}_{2^{-m}}(E_r)=\mathcal{M}^{s}_{2^{-r}}(E_r)=\mathcal{M}^{s}_{2^{-m}}(E_m)\qquad(r\geq m).
\end{equation}
令 $F$ 为紧集 $F=\bigcap_{m}^{\infty}E_k$.于是利用\cref{thm:5.1} 与 \eqref{eq:5.10},
\begin{equation}\label{eq:5.12}
  \mathcal{H}^{s}(F)\leq\mathcal{M}^{s}(F)=\lim_{k\to\infty}\mathcal{M}^{s}_{2^{-k}}(F)
  \leq\varliminf_{k\to\infty}\mathcal{M}^{s}_{2^{-k}}(E_k)=\mathcal{M}^{s}_{2^{-m}}(E_m)<\infty.
\end{equation}
另一方面，利用\cref{lem:5.2}、\cref{thm:5.1} 与 \eqref{eq:5.11},
\[
  2^{s}\mathcal{H}^{s}(F)\geq 2^{s}\mathcal{H}^{s}_{2^{-(m+1)}}(F)\geq\lim_{k\to\infty}\mathcal{H}^{s}_{2^{-m}}(E_k)
  \geq b_1^{-1}\lim_{k\to\infty}\mathcal{M}^{s}_{2^{-m}}(E_k)
  =b_1^{-1}\mathcal{M}^{s}_{2^{-m}}(E_m)\geq 2^{s}c.
\]
合并这些不等式得
\[
  c\leq\mathcal{H}^{s}(F)<\infty.
\]
若左端是严格不等式，则利用测度 $\mathcal{H}^s$ 的上、下连续性（\cref{thm:1.1}），形如 $F\cap[u,\infty)$ 的集合将有测度恰好为 $c$.

上述方法的一个微小变体证明 (b).设 $J$ 是长度为 $2^{-r}$ 的闭二进区间，$r\geq m$.若把 \eqref{eq:5.9} 对位于 $J$ 内的长度为 $2^{-k}$ 的二进子区间 $I$ 求和（$k\geq r$）并照前进行，则在 \eqref{eq:5.10} 的位置得到
\[
  \mathcal{M}^{s}_{2^{-k}}(E_k\cap J)=\mathcal{M}^{s}_{2^{-m}}(E_m\cap J)\qquad(k\geq r\geq m),
\]
从而在 \eqref{eq:5.12} 的位置得到
\[
  \mathcal{H}^{s}(F\cap J)\leq\mathcal{M}^{s}_{2^{-r}}(E_r\cap J)\leq\mathcal{M}^{s}_{2^{-r}}(J)\leq|J|^{s}.
\]
任何区间 $J_0$ 都可被至多六个长度至多为 $|J_0|$ 的相邻二进区间覆盖，由此即得结论.(b) 的完整细节留给读者.
\end{proof}

鉴于\cref{thm:5.4}，一个显然要问的问题是：这类结果对更一般的集合类是否成立.Besicovitch 在其原始文章中把该定理推广到了 $F_{\sigma\delta\sigma}$-集.\textcite{davies1952b} 利用下述中间结果把这一工作推进到其自然的终点，即对一切 Souslin 集证明该定理.

\begin{theorem}\label{thm:5.5}
设 $E$ 是 $\mathbb{R}^n$ 中满足 $\mathcal{H}^s(E)=\infty$ 的任意 Souslin 子集.则 $E$ 包含一个 $\mathcal{H}^s$-测度无穷的闭子集.
\end{theorem}

\begin{proof}
见 \textcite{davies1952b} 或 \textcite{rogers1970}.
\end{proof}

与\cref{thm:5.4} 结合可得：

\begin{theorem}\label{thm:5.6}
设 $E$ 是 $\mathbb{R}^n$ 中满足 $\mathcal{H}^s(E)=\infty$ 的 Souslin 子集.

(a) 对任意正常数 $c$，存在 $E$ 的紧子集 $F$ 使 $\mathcal{H}^s(F)=c$.

(b) 存在 $E$ 的紧子集 $F$ 使 $\mathcal{H}^s(F)>0$，且对某个常数 $b$ 有
\[
  \mathcal{H}^{s}\bigl(B_r(x)\cap F\bigr)\leq br^{s}\qquad(x\in\mathbb{R}^n,\ r\leq 1).
\]
\end{theorem}

本节最后提到一个更具病态性质的结果：若 $E$ 是 $\mathcal{H}^s$-测度非 $\sigma$-有限的紧集，则 $E$ 大到可以分解为连续统多个 $\mathcal{H}^s$-测度非 $\sigma$-有限的子集.\textcite{davies1968} 与 \textcite{rogers1970}（2.8 节）提供了这一思想及相关思想的进一步细节.

\section{集合的 Descartes 积}

为记号简便，本节始终把注意力限于平面的子集，尽管这些工作可以毫无困难地推广到更高维.这里 $(x,y)$ 表示 $\mathbb{R}^2$ 中的 Descartes 坐标.若 $E$ 是 $\mathbb{R}^2$ 的任意子集，用 $E_x$ 表示由 $E$ 中第一坐标等于 $x$ 的那些点构成的截口.

若 $E$ 是 $\mathbb{R}^2$ 的勒贝格可测子集，\newterm{fubinis-theorem}{Fubini 定理}{Fubini's theorem}的一个直接推论是：$E_x$ 对几乎所有的 $x$ 是 $\mathcal{L}^1$-可测的，且
\[
  \mathcal{L}^{2}(E)=\int\mathcal{L}^{1}(E_x)\,\mathrm{d}\mathscr{L}^{1}(x),
\]
其中 $\mathscr{L}^1$ 与 $\mathcal{L}^2$ 分别是 $1$ 维与 $2$ 维勒贝格测度.自然会问：把勒贝格测度换成分数维数的 Hausdorff 测度后，是否有任何类似的结果.例如，若 $A$ 与 $B$ 是两条坐标轴的子集，$\mathcal{H}^{s+t}(A\times B)$ 与 $\mathcal{H}^{s}(A)\mathcal{H}^{t}(B)$ 是否相等，至少相差一个常数倍？事实上，正如后面的例子所示，这类结果远非成立.一种初等的证明尝试或许是：用长度为 $\delta$ 的区间覆盖 $A$ 与 $B$，诱导出 $A\times B$ 的由边长 $\delta$ 的正方形构成的覆盖，从而估计各个 Hausdorff 测度.遗憾的是，使 $A$ 的覆盖``有效''的区间长度可能使 $B$ 的覆盖``无效''，使这类估计毫无用处.（当然，乘积测度 $\mathcal{H}^s\times\mathcal{H}^t$ 未必是 Hausdorff 测度.）不过，得到如下形式的不等式是可能的：
\[
  \mathcal{H}^{s+t}(A\times B)\geq b\,\mathcal{H}^{s}(A)\mathcal{H}^{t}(B).
\]
\textcite{besicovitchmoran1945} 在一定条件下证明了这一点；其他基于密度思想的早期证明由 Moran（\citeyear{moran1946,moran1949}）、\textcite{freilich1950}、\textcite{eggleston1950a} 与 \textcite{eggleston1953b} 给出；另见 \textcite{eggleston1950b}.然而这些证明相当技术化，而且对 $A$ 与 $B$ 附加了限制.下面这个使用可比较网测度的非常一般而漂亮的方法归功于 \textcite{marstrand1954b}.

证明依赖于下述组合引理.

\begin{lemma}\label{lem:5.7}
设 $A$ 是 $\mathbb{R}$ 的任意子集，$\{I_i\}$ 是 $A$ 的由二进区间构成的可数 $\delta$-覆盖，$\{a_i\}$ 是正数序列.设 $c$ 是使
\begin{equation}\label{eq:5.13}
  \sum_{\{i:x\in I_i\}}a_i>c
\end{equation}
对一切 $x\in A$ 成立的常数.则
\begin{equation}\label{eq:5.14}
  \sum_{i}a_i|I_i|^{s}\geq c\,\mathcal{M}^{s}_{\delta}(A).
\end{equation}
\end{lemma}

\begin{proof}
先设族 $\{I_i\}$ 与 $\{a_i\}$ 有限.把 $a_i$ 稍微减小，可取每个 $a_i$ 为有理数而不失去 \eqref{eq:5.13}；再乘以一个通分因子，可进一步设 $a_i$ 为整数.于是对每个 $i$ 取 $a_i$ 份区间 $I_i$，便可设 $a_i=1$ 对一切 $i$ 成立.

在这些假设下，\eqref{eq:5.13} 蕴含 $A$ 中每一点 $x$ 至少属于 $\lceil c\rceil$ 个区间，其中 $\lceil c\rceil$ 是大于 $c$ 的最小整数.由于 $A\subset\bigcup_i I_i$ 且 $\{I_i\}$ 是具有网性质的二进区间，可以取所有不被族中其他区间包含的区间，得到包含 $A$ 的互不重叠的区间组，记作 $\{I_i\}_{i\in\mathcal{S}_1}$.剩下的区间仍至少 $\lceil c\rceil-1$ 次覆盖 $A$，故可用同样的方法选出它们的不重叠子族 $\{I_i\}_{i\in\mathcal{S}_2}$，它也覆盖 $A$.如此进行，得到 $\lceil c\rceil$ 个互不重叠的区间族 $\{I_i\}_{i\in\mathcal{S}_j}$（$j=1,2,\ldots,\lceil c\rceil$），每个族都覆盖 $A$，且指标集 $\mathcal{S}_j$ 两两不交.由于对一切 $i$ 有 $|I_i|\leq\delta$，
\[
  \sum_{i\in\mathcal{S}_j}|I_i|^{s}\geq\mathcal{M}^{s}_{\delta}(A),
\]
对所有 $j$ 求和，便得到 $\{I_i\}$ 为有限区间族情形的 \eqref{eq:5.14}.

最后，设 $\{I_i\}_{1}^{\infty}$ 是无穷的二进区间族.对每个 $k$ 令
\[
  A_k=\Bigl\{x\in\mathbb{R}:\sum_{\substack{i:\,x\in I_i\\ i\leq k}}a_i>c\Bigr\}.
\]
由有限情形，
\[
  \sum_{1}^{k}a_i|I_i|^{s}\geq c\,\mathcal{M}^{s}_{\delta}(A_k).
\]
但每个 $A_k$ 都是二进区间的有限并，序列 $\{A_k\}$ 递增，且 $A\subset\bigcup_k A_k=\lim A_k$.由\cref{lem:5.3},
\[
  \sum_{1}^{\infty}a_i|I_i|^{s}
  \geq c\lim_{k\to\infty}\mathcal{M}^{s}_{\delta}(A_k)
  =c\,\mathcal{M}^{s}_{\delta}\Bigl(\lim_{k\to\infty}A_k\Bigr)
  \geq c\,\mathcal{M}^{s}_{\delta}(A),
\]
无穷情形证毕.
\end{proof}

\begin{theorem}\label{thm:5.8}
设 $E$ 是平面集合，$A$ 是 $x$ 轴的任意子集.设若 $x\in A$ 则 $\mathcal{H}^1(E_x)>c$，$c$ 为某常数.则
\[
  \mathcal{H}^{s+t}(E)\geq bc\,\mathcal{H}^{s}(A),
\]
其中 $b$ 只依赖于 $s$ 与 $t$.
\end{theorem}

\begin{proof}
由\cref{thm:5.1}，把结果中的 `H' 全部换成 `M' 来证明即可.

给定 $\delta>0$，设 $\{S_i\}_{1}^{\infty}$ 是构成 $E$ 的 $2^{\frac{1}{2}}\delta$-覆盖的二进正方形族.对每个 $x\in A$，
\[
  E_x\subset\bigcup_{1}^{\infty}(S_i)_x,
\]
故
\[
  \mathcal{M}^{t}_{\delta}(E_x)\leq\sum_{1}^{\infty}|(S_i)_x|^{t}.
\]
若 $A_{\delta}=\{x\in A:\mathcal{M}^{t}_{\delta}(E_x)>c\}$，则对 $x\in A_{\delta}$，
\[
  c<\sum_{1}^{\infty}|(S_i)_x|^{t}
  =2^{-\frac{1}{2}t}\sum_{\{i:\,x\in\operatorname{proj}S_i\}}|S_i|^{t},
\]
其中 $\operatorname{proj}S_i$ 是把 $S_i$ 投影到 $x$ 轴上所得的长度至多为 $\delta$ 的二进区间.于是
\[
  \sum_{1}^{\infty}|S_i|^{s+t}
  =\sum_{1}^{\infty}|S_i|^{t}|S_i|^{s}
  =2^{\frac{1}{2}s}\sum_{1}^{\infty}|S_i|^{t}|\operatorname{proj}S_i|^{s}
  \geq 2^{\frac{1}{2}(s+t)}c\,\mathcal{M}^{s}_{\delta}(A_{\delta}),
\]
其中在\cref{lem:5.7} 中取 $I_i=\operatorname{proj}S_i$，$a_i=|S_i|^t$.这对 $E$ 的任何由二进正方形 $\{S_i\}$ 构成的 $2^{\frac{1}{2}}\delta$-覆盖都成立，故
\[
  bc\,\mathcal{M}^{s}_{\delta}(A_{\delta})\leq\mathcal{M}^{s+t}_{2^{\frac{1}{2}}\delta}(E)\leq\mathcal{M}^{s+t}(E),
\]
其中 $b=2^{\frac{1}{2}(s+t)}$.由于当 $\delta$ 递减到 $0$ 时 $A_{\delta}$ 递增到 $A$，故当 $\delta\leq\rho$ 时
\[
  \mathcal{M}^{s}_{\delta}(A_{\rho})\leq\mathcal{M}^{s}_{\delta}(A_{\delta})\leq b^{-1}c^{-1}\mathcal{M}^{s+t}(E).
\]
于是对 $\rho>0$，
\[
  \mathcal{M}^{s}(A_{\rho})\leq b^{-1}c^{-1}\mathcal{M}^{s+t}(E),
\]
再由测度 $\mathcal{M}^s$ 的连续性，
\[
  \mathcal{M}^{s}(A)\leq b^{-1}c^{-1}\mathcal{M}^{s+t}(E). \qedhere
\]
\end{proof}

现在取 $E$ 为\newterm{cartesian-product}{Descartes 积}{Cartesian product} $A\times B$.

\begin{corollary}\label{cor:5.9}
对 $\mathbb{R}$ 的任意子集 $A$ 与 $B$，
\[
  \mathcal{H}^{s+t}(A\times B)\geq b\,\mathcal{H}^{s}(A)\mathcal{H}^{t}(B).
\]
\end{corollary}

可以用 Hausdorff 维数来解释这一结果：

\begin{corollary}\label{cor:5.10}
对 $\mathbb{R}$ 的任意子集 $A$ 与 $B$，
\[
  \dim(A\times B)\geq\dim A+\dim B.
\]
\end{corollary}

注意，当 $t=0$（此时 $\mathcal{H}^t(E_x)$ 等于 $E_x$ 中点的个数）时这些结果也成立.常数 $b$ 的值由 \textcite{ernstfreilich1976} 讨论.

一般而言，\cref{thm:5.8} 及其推论中的不等式不可反向，尽管 \textcite{besicovitchmoran1945} 与 \textcite{taylor1952} 给出了使\cref{cor:5.10} 中等号成立的关于 $A,B$ 的若干充分条件.下述关于该不等式可以严格的演示，作为 \textcite{besicovitchmoran1945} 例子的一种简单替代.

\begin{theorem}\label{thm:5.11}
存在 Hausdorff 维数为 $0$ 的 Borel 子集 $A,B\subset\mathbb{R}$，使 $\mathcal{H}^1(A\times B)>0$.
\end{theorem}

\begin{proof}
设 $\{s_j\}$ 是递减于 $0$ 的数列，并设 $0=m_0<m_1<m_2<\cdots$ 是增长足够快以保证
\begin{equation}\label{eq:5.15}
  \left.
  \begin{aligned}
    &(m_1-m_0)+(m_3-m_2)+\cdots+(m_{2j-1}-m_{2j-2})\leq s_jm_{2j}\\
    &(m_2-m_1)+(m_4-m_3)+\cdots+(m_{2j}-m_{2j-1})\leq s_jm_{2j+1}
  \end{aligned}
  \right\}
\end{equation}
的整数列.设 $A$ 为 $[0,1]$ 中如下数构成的子集：当 $m_j+1\leq r\leq m_{j+1}$ 且 $j$ 为奇数时，其第 $r$ 位小数为零.类似地，取 $B$ 为当 $m_j+1\leq r\leq m_{j+1}$ 且 $j$ 为偶数时第 $r$ 位小数为零的数之集.取显然的覆盖：用 $10^k$ 个长度为 $10^{-m_{2j}}$ 的区间覆盖 $A$，其中
\[
  k=(m_1-m_0)+(m_3-m_2)+\cdots+(m_{2j-1}-m_{2j-2}),
\]
则由 \eqref{eq:5.15}，当 $s>0$ 时 $\mathcal{H}^s(A)=0$；类似地 $\mathcal{H}^s(B)=0$.

设 $\operatorname{proj}$ 表示从平面到直线 $L$（即 $y=x$）的正交投影.则 $\operatorname{proj}(x,y)$ 是 $L$ 上与原点的（带符号）距离为 $2^{-\frac{1}{2}}(x+y)$ 的点.若 $u\in[0,1]$，则可找到 $x\in A$ 与 $y\in B$ 使 $u=x+y$（$u$ 的某些小数位由 $x$ 提供，其余由 $y$ 提供）.于是 $\operatorname{proj}(A\times B)$ 是 $L$ 的长度为 $2^{-\frac{1}{2}}$ 的子区间.利用正交投影不增大距离、从而由\cref{lem:1.8} 也不增大 Hausdorff 测度这一事实，
\[
  2^{-\frac{1}{2}}=\mathcal{H}^{1}\bigl(\operatorname{proj}(A\times B)\bigr)\leq\mathcal{H}^{1}(A\times B).
\]
如果需要，可以通过添加可数点集把 $A$ 与 $B$ 变成紧集.
\end{proof}

本节的全部结果都可以在更一般的框架中陈述；见 \textcite{larman1967a} 与 Wegmann（\citeyear{wegmann1969a,wegmann1969b,wegmann1971a}），他们考虑了非常一般的空间中集合的乘积.特别地，\cref{cor:5.9} 的高维类似物可用类似方法证明，即：

\begin{theorem}\label{thm:5.12}
设 $A\subset\mathbb{R}^n$，$B\subset\mathbb{R}^m$.则对某个常数 $b$ 有
\[
  \mathcal{H}^{s+t}(A\times B)\geq b\,\mathcal{H}^{s}(A)\mathcal{H}^{t}(B),
\]
其中 $A\times B\subset\mathbb{R}^{n+m}$.
\end{theorem}

\newterm{cylinder-set}{柱集}{cylinder set}是 $\mathbb{R}^3$ 中形如 $E\times I$ 的子集，其中 $E$ 是 $\mathbb{R}^2$ 的 $1$ 维子集，$I\subset\mathbb{R}$ 是单位区间.\textcite{randolph1936} 提出的一个重要问题是确定何时
\[
  \mathcal{H}^{1}(E)=\frac{1}{4}\pi\,\mathcal{H}^{2}(E\times I),
\]
即何时``面积 = 底 $\times$ 高''.有兴趣的读者可参阅 \textcites{besicovitchmoran1945,freilich1965,ward1967,larman1967c} 了解进一步细节.

\begin{exercises}
\item 设 $E\subset\mathbb{R}$ 是 $s$-集，$I$ 是单位区间.证明 $E\times I$ 是 $\mathbb{R}^2$ 中的 $(s+1)$-集.
\item 记号如\cref{thm:5.8} 所设，证明若 $E$ 是 $\mathbb{R}^2$ 中的 Borel 集，则
\[
  \mathcal{H}^{s+t}(E)\geq b\int\mathcal{H}^{t}(E_x)\,\mathrm{d}\mathcal{H}^{s}(x).
\]
\item 设 $E\subset\mathbb{R}$ 是由 Cantor 型过程构成的 $s$-集，$E=\bigcap_{1}^{\infty}E_j$，其中每个 $E_j$ 有等长区间构成的 $\delta$-覆盖 $\{U_i\}$ 使 $\sum|U_i|^s\leq c$，且当 $j\to\infty$ 时 $\delta\to 0$.设 $F=\{x-y:x,y\in E\}$ 为 $E$ 的差集.证明 $\dim F\leq 2\dim E$.（提示：考虑 $E\times E$ 到直线 $x+y=0$ 上的投影.）
\item 由\cref{exe:1.7} 推出：集合
\[
  G=\{(x,y):y-x\ \text{是无理数}\}\subset\mathbb{R}^2
\]
具有全平面勒贝格测度，但不含任何形如 $A\times B$ 的子集，其中 $A$ 与 $B$ 都是正测度的勒贝格可测集.
\end{exercises}
